% Generated draft: requires mathematical and source review.
\documentclass[11pt]{article}
\usepackage[margin=25mm]{geometry}
\usepackage{fontspec}
\setmainfont{Latin Modern Roman}
\usepackage{amsmath,amssymb,mathtools,unicode-math}
\setmathfont{Latin Modern Math}
\usepackage{enumitem,needspace,xurl,hyperref}
\setlength{\parindent}{0pt}
\setlength{\parskip}{5pt}
\setlength{\emergencystretch}{3em}
\newcommand{\N}{\mathbb{N}}
\newcommand{\Z}{\mathbb{Z}}
\newcommand{\Q}{\mathbb{Q}}
\newcommand{\R}{\mathbb{R}}
\newcommand{\C}{\mathbb{C}}
\newcommand{\F}{\mathbb{F}}
\newcommand{\A}{\mathbb{A}}
\newcommand{\PP}{\mathbb{P}}
\newcommand{\E}{\mathbb{E}}
\newcommand{\eps}{\varepsilon}
\newcommand{\dd}{\,\mathrm{d}}
\begin{document}
\section*{UnsolvedMath ID: 5200009}
\textbf{Source problem code:} catalogue entry
\subsection{Definitions and mathematical statement}
Let $\varepsilon>0$, let $\Lambda_\varepsilon=\varepsilon\Z^2$ be the square lattice, and let $C\subset\R^2$ be a curve. For a prescribed translation path $a:I\to\R^2$, the click-event set is
\[
 E(C,a,\varepsilon)=\{t\in I:(C+a(t))\cap\Lambda_\varepsilon\ne\varnothing\}.
\]
One may also record the multiplicity $\#((C+a(t))\cap\Lambda_\varepsilon)$ when finite. The source does not specify the translation path, speed, parameter interval, or whether simultaneous contacts produce one click or several. Straight translation $a(t)=tv$ is a possible special case, not an assumption supplied by the record. Without curve and genericity hypotheses, click events need not form a locally finite sequence.

The source asks whether the shape of $C$ can be recovered from the distribution of click parameters, either using $C$ alone or using all its rotations. If $R_\theta$ is planar rotation by $\theta$, the latter would involve data from $E(R_\theta C,a,\varepsilon)$ for all intended angles $\theta$. It is not stated whether the angle labels are part of the data. Nor is ``distribution'' specified: full event times, spacing statistics, an empirical measure, or a limiting probability law contain different information.

The recovery question therefore asks for a suitable data map $\mathcal D$ and curve class for which
\[
 \mathcal D(C_1)=\mathcal D(C_2)
 \quad\Longrightarrow\quad
 C_1\text{ and }C_2\text{ have the same intended shape},
\]
with the equivalence of shapes and the translation/rotation data to be fixed. The source separately asks for a Fourier-type transform extracting curvature patterns from the click distribution. Its transform, curvature notion, and recovery or detection criterion are not prescribed. Finally it explicitly requests analysis of the segment, polygon, and circle cases; these are retained as distinct test cases, not replaced by a smooth-curve-only question.
\subsection{Short English statement}
Can lattice-contact data from translating a curve reveal its shape, or reveal its curvature through a Fourier-like transform? Study data from one orientation and from all rotations, including segments, polygons, and circles.

The supplied research summary describes the problem as exploratory ``big data'' questions and reports finding no formal recovery or Fourier-extraction theorem. It mentions segment intuition involving rotations/Brill sequences and circle intuition involving Gauss circle counting, without exact results or citations. The OPEN-TRIAGE label and original 2021 open status are supplied assessments; provenance records no later authoritative resolution check.
\subsection{Sources}
\begin{itemize}
\item Local UnsolvedMath registry, permanent ID 5200009, source code catalogue entry, attributed to Mark Levi, 2021. Dataset curation and metadata: Ulam AI and UnsolvedMath Contributors, CC BY 4.0.
\item Bialy, Fierobe, Glutsyuk, Levi, Plakhov, and Tabachnikov, \emph{Open problems on billiards and geometric optics} (2021), Mark Levi, Problem 1, source-order ordinal 9, \url{https://arxiv.org/abs/2110.10750}; source-TeX extraction, accessed 2026-07-29.
\item The supplied summary's rotation/Brill-sequence and Gauss-circle-counting associations have no separate primary citations or precise statements in the record.
\item Source provenance marks status evidence NEEDS\_REVIEW and rights NEEDS\_REVIEW: the source presents the item as open, no later authoritative resolution check is recorded, and redistribution permission was not established despite verified public access. This qualification is separate from the dataset licence.
\end{itemize}
\end{document}
